\documentclass[12pt,a4paper]{article}
\usepackage{amsmath,amssymb,amsthm}

\begin{document}
	
	\section*{Cyclic Groups: Formula Sheet}
	
	\subsection*{Definition}
	A \emph{cyclic group} is a group $G$ generated by a single element $g \in G$.  
	\[
	G = \langle g \rangle =
	\begin{cases}
		\{ g^n \mid n \in \mathbb{Z} \}, & \text{(multiplicative notation)} \\[6pt]
		\{ n g \mid n \in \mathbb{Z} \}, & \text{(additive notation).}
	\end{cases}
	\]
	
	\subsection*{Identity and Inverses}
	- Identity element:
	\[
	g^0 = e \quad \text{(multiplicative)}, \qquad 0g = e \quad \text{(additive)}.
	\]
	
	- Inverses:
	\[
	(g^n)^{-1} = g^{-n}, \qquad (ng)^{-1} = (-n)g.
	\]
	
	\subsection*{Order of a Cyclic Group}
	- The \emph{order} of $\langle g \rangle$ is the number of distinct elements generated by $g$.  
	
	- If no finite order exists, then
	\[
	\langle g \rangle \cong (\mathbb{Z}, +).
	\]
	
	- If $g$ has finite order $n$, then
	\[
	\langle g \rangle \cong (\mathbb{Z}_n, +).
	\]
	
	\subsection*{Criterion for Generators}
	If $G = \langle g \rangle$ has order $n$, then
	\[
	g^k \text{ is a generator of } G \quad \Longleftrightarrow \quad \gcd(k,n) = 1.
	\]
	
	\subsection*{Subgroup Structure}
	Every subgroup of a cyclic group is cyclic.  
	If $|G| = n$, then for each divisor $d \mid n$ there is a unique subgroup of order $d$:
	\[
	H_d = \langle g^{n/d} \rangle.
	\]
	
	\subsection*{Examples}
	
	\paragraph{1. Integers modulo $6$:}
	\[
	(\mathbb{Z}_6, +), \qquad 
	\langle 1 \rangle = \{0,1,2,3,4,5\}.
	\]
	Other subgroups:
	\[
	\langle 2 \rangle = \{0,2,4\}, \qquad 
	\langle 3 \rangle = \{0,3\}.
	\]
	
	\paragraph{2. $6$th roots of unity:}
	\[
	\langle \zeta \rangle = \{ \zeta^k \mid k=0,1,2,3,4,5 \}, \qquad 
	\zeta = e^{2\pi i/6}.
	\]
	This is a cyclic group of order $6$ under multiplication in $\mathbb{C}$.
	
	\subsection*{Key Structural Result}
	Every finite abelian group can be expressed as a direct product of cyclic groups.  
	Thus cyclic groups are the fundamental building blocks of abelian group theory.
	
\end{document}
